\subsection{Non-ambient laws over every \texorpdfstring{$K$}{K}}\label{dp:ssec:degree-p}

For the proof of Proposition~\ref{dp:prop:degree-p} we extend the notation of Section~\ref{dp:ssec:tensor-notation} from $\Bbbk^m$ to $\Bbbk^{\mathcal J}$, for a field $\Bbbk$ and a finite set $\mathcal J$. A tensor $t\in(\Bbbk^{\mathcal J})^{\otimes d}$ is read as a function on the words $w=(w_1,\dots,w_d)\in\mathcal J^d$. The space $\Gamma^d(\Bbbk^{\mathcal J})$ is the space of symmetric tensors, and $\Gamma^0(\Bbbk^{\mathcal J})=\Bbbk$. Let $\N^{\mathcal J}_d$ be the set of $\alpha\in\N^{\mathcal J}$ with $\abs\alpha=d$; for $\mathcal J=[m]$ it is the set $\N^m_d$ of Section~\ref{dp:sec:nonambient}. The \emph{content} of a word $w\in\mathcal J^d$ is the $\alpha\in\N^{\mathcal J}_d$ whose entry at each element of $\mathcal J$ is the number of letters of $w$ equal to it. Let $O_\alpha\in\Gamma^d(\Bbbk^{\mathcal J})$ be the sum, over the words $w$ of content $\alpha$, of the tensor product of the unit vectors of $\Bbbk^{\mathcal J}$ at the letters $w_1,\dots,w_d$. The $O_\alpha$, $\alpha\in\N^{\mathcal J}_d$, form a basis of $\Gamma^d(\Bbbk^{\mathcal J})$, and $t=\sum_\alpha t_\alpha O_\alpha$, where $t_\alpha$ is the value of $t$ at any word of content $\alpha$.

The \emph{product} of $t\in\Gamma^d(\Bbbk^{\mathcal J})$ and $t'\in\Gamma^{d'}(\Bbbk^{\mathcal J})$ is the sum $t\cdot t'$ of the $\binom{d+d'}d$ tensors obtained from $t\otimes t'$ by the permutations of the $d+d'$ factors that keep the order of the first $d$ factors and the order of the last $d'$. A scalar in $\Gamma^0(\Bbbk^{\mathcal J})$ multiplies. These permutations represent the cosets of the subgroup that permutes the first $d$ factors and the last $d'$ factors among themselves, and this subgroup fixes $t\otimes t'$. So $t\cdot t'$ is symmetric, and it lies in $\Gamma^{d+d'}(\Bbbk^{\mathcal J})$. For $\Bbbk=K$ this is the shuffle product of Section~\ref{dp:sec:conv}. Tensor powers of linear maps commute with permutations of the factors, so $q^{\otimes(d+d')}(t\cdot t')=q^{\otimes d}(t)\cdot q^{\otimes d'}(t')$ for every linear map $q$ of $\Bbbk^{\mathcal J}$ into some $\Bbbk^{\mathcal J'}$.

Let $\mathcal J=[m]$, $\alpha\in\N^m_d$ and $\alpha'\in\N^m_{d'}$. The value of $O_\alpha\cdot O_{\alpha'}$ at a word of content other than $\alpha+\alpha'$ is $0$. Its value at a word of content $\alpha+\alpha'$ is the number of ways to choose the $d$ of its letters that come from the first factor, which must contain $\alpha_i$ of the $\alpha_i+\alpha'_i$ letters equal to $i$, for each $i$. So
\begin{equation}\label{dp:eq:orbit-product}
 O_\alpha\cdot O_{\alpha'}=\prod_{i=1}^m\binom{\alpha_i+\alpha'_i}{\alpha_i}\,O_{\alpha+\alpha'},
\end{equation}
and the product is commutative and associative on the spaces $\Gamma^d(\Bbbk^m)$, $d\ge0$, as it is bilinear and $\binom{a+b}a\binom{a+b+c}{a+b}=(a+b+c)!/(a!\,b!\,c!)$ for integers $a,b,c\ge0$. For $t\in\Gamma^{d+d'}(\Bbbk^m)$ and $\gamma\in\N^m_d$ let $\partial^\gamma t\in\Gamma^{d'}(\Bbbk^m)$ have the coefficients $(\partial^\gamma t)_\alpha=t_{\alpha+\gamma}$, $\alpha\in\N^m_{d'}$. For $p=\operatorname{char}K>0$, $\Bbbk=\F_p$, $d=p-1$ and $d'=1$ this is the slice of Section~\ref{dp:sec:nonambient}. Then $\partial^\gamma\partial^{\gamma'}=\partial^{\gamma+\gamma'}$. Let $\alpha\in\N^m_{d+d'}$. Then $\partial^\gamma O_\alpha$ is $O_{\alpha-\gamma}$ if $\gamma\le\alpha$ entrywise and $0$ otherwise. By~\eqref{dp:eq:orbit-product}, $(\partial^\gamma O_\alpha)\cdot O_\gamma=\prod_i\binom{\alpha_i}{\gamma_i}O_\alpha$ for $\gamma\le\alpha$. Vandermonde's identity gives $\sum_\gamma\prod_i\binom{\alpha_i}{\gamma_i}=\binom{d+d'}d$, the sum over the $\gamma\in\N^m_d$ with $\gamma\le\alpha$. Hence, by linearity in $t$,
\begin{equation}\label{dp:eq:contraction-sum}
 \sum_{\gamma\in\N^m_d}(\partial^\gamma t)\cdot O_\gamma=\binom{d+d'}d\,t\qquad(t\in\Gamma^{d+d'}(\Bbbk^m)).
\end{equation}

\begin{lemma}[an identity for symmetric tensors]\label{dp:lem:linear}
Let $\Bbbk$ be a field of characteristic $p>0$ and $m\ge1$, let $\mathsf L$ be a matrix over $\Bbbk$ with rows and columns indexed by $\N^m_{p-1}$, and let $\mathsf L^\sharp$ be the linear map of $\Gamma^p(\Bbbk^m)$ with
\[
 \mathsf L^\sharp(t)=\sum_{\gamma,\gamma'\in\N^m_{p-1}}\mathsf L_{\gamma\gamma'}\,(\partial^\gamma t)\cdot O_{\gamma'}.
\]
\begin{enumerate}
\item We have $t=\sum_{\gamma\in\N^m_{p-1}}\bigl((\id-\mathsf L^\sharp)(\partial^\gamma t)\bigr)\cdot O_\gamma$ for every $t\in\Gamma^{2p-1}(\Bbbk^m)$.
\item We have $\rank(\id-\mathsf L^\sharp)\ge\binom{m+2p-2}{2p-1}\big/\binom{m+p-2}{p-1}\ge(p-1)!\,m^p/(2p-1)!$.
\end{enumerate}
\end{lemma}

\begin{proof}
\emph{Step 1 (the identity part).} Let $t\in\Gamma^{2p-1}(\Bbbk^m)$. By~\eqref{dp:eq:contraction-sum} with $d=p-1$ and $d'=p$, we have $\sum_\gamma(\partial^\gamma t)\cdot O_\gamma=\binom{2p-1}{p-1}t$, the sum over $\N^m_{p-1}$. The integer $\binom{2p-1}{p-1}(p-1)!$ is the product $(p+1)(p+2)\cdots(2p-1)$, which is congruent to $(p-1)!$ modulo $p$, and $(p-1)!$ is prime to $p$. So $\binom{2p-1}{p-1}=1$ in $\Bbbk$, and $\sum_\gamma(\partial^\gamma t)\cdot O_\gamma=t$.

\emph{Step 2 (the part of $\mathsf L^\sharp$).} For $\gamma\in\N^m_{p-1}$ we have $\mathsf L^\sharp(\partial^\gamma t)=\sum_{\gamma',\gamma''}\mathsf L_{\gamma'\gamma''}\,(\partial^{\gamma'}\partial^\gamma t)\cdot O_{\gamma''}$. The product is commutative and associative, and $\partial^{\gamma'}\partial^\gamma=\partial^\gamma\partial^{\gamma'}$. So
\[
 \sum_\gamma\mathsf L^\sharp(\partial^\gamma t)\cdot O_\gamma
 =\sum_{\gamma',\gamma''}\mathsf L_{\gamma'\gamma''}\,O_{\gamma''}\cdot
 \sum_\gamma\bigl(\partial^\gamma(\partial^{\gamma'}t)\bigr)\cdot O_\gamma,
\]
all sums over $\N^m_{p-1}$. By~\eqref{dp:eq:contraction-sum} with $d=p-1$ and $d'=1$, applied to $\partial^{\gamma'}t\in\Gamma^p(\Bbbk^m)$, the inner sum is $p\,\partial^{\gamma'}t=0$. With Step~1 this is item~(1).

\emph{Step 3 (item (2)).} By item~(1), the linear map $t\mapsto\bigl((\id-\mathsf L^\sharp)(\partial^\gamma t)\bigr)_{\gamma\in\N^m_{p-1}}$ from $\Gamma^{2p-1}(\Bbbk^m)$ into the product of $\abs{\N^m_{p-1}}$ copies of the image of $\id-\mathsf L^\sharp$ is injective. As $\dim\Gamma^d(\Bbbk^m)=\abs{\N^m_d}=\binom{m+d-1}d$, this gives $\binom{m+2p-2}{2p-1}\le\binom{m+p-2}{p-1}\rank(\id-\mathsf L^\sharp)$. The quotient of the two binomial coefficients is $(p-1)!/(2p-1)!$ times the product of the $p$ numbers $m+p-1,\dots,m+2p-2$, so it is at least $(p-1)!\,m^p/(2p-1)!$.
\end{proof}

\begin{proof}[Proof of Proposition~\ref{dp:prop:degree-p}]
Let $p=\operatorname{char}K>0$ and $m\ge1$, and let $\mathcal J_m$ be the disjoint union of $[m]$ and $\N^m_{p-1}$. For a field $\Bbbk$ containing $\F_p$ write $e_i$, $i\in[m]$, and $e_\gamma$, $\gamma\in\N^m_{p-1}$, for the unit vectors of $\Bbbk^{\mathcal J_m}$, and read $\Bbbk^m$ as the span of the $e_i$. For $T\in\Gamma^p(\F_p^m)$ let $q_T\colon\Bbbk^{\mathcal J_m}\to\Bbbk^m$ be the linear map with $q_T(e_i)=e_i$ and $q_T(e_\gamma)=\partial^\gamma T$. Its matrix has entries in $\F_p$. For a real number $\tau\ge1$ and $z=\sum_iz_ie_i+\sum_\gamma z_\gamma e_\gamma$ in $K^{\mathcal J_m}$ put $\norm z_\tau=\sum_i\abs{z_i}+\tau\sum_\gamma\abs{z_\gamma}$. By~\eqref{dp:eq:weighted-norm}, $\norm h_{T,\tau}=\inf\{\norm z_\tau:q_T(z)=h\}$.

Call the elements $\mathbf a$ of $\N^{\mathcal J_m}_p$ \emph{positions}. Let $\abs{\mathbf a}_\partial$ be the sum of the entries of $\mathbf a$ at the elements of $\N^m_{p-1}$, and call $\mathbf a$ \emph{nonlinear} if $\abs{\mathbf a}_\partial\ge2$. For a word $w\in\mathcal J_m^p$ let $\abs w_\partial$ be the number of its letters in $\N^m_{p-1}$. For $S\in\Gamma^p(\Bbbk^{\mathcal J_m})$ let $\operatorname{supp}S$ be the set of positions $\mathbf a$ with $S_{\mathbf a}\ne0$. Put $\mathsf r_m=(p-1)!\,m^p/(2p-1)!$. In Steps~1 to~3, $T\in\Gamma^p(\F_p^m)$ and $\tau\ge1$ are arbitrary.

\emph{Step 1 (norms).} If $h=q_T(z)$, the $i$-th entry of $h$ is $z_i+\sum_\gamma z_\gamma T_{\gamma^{+i}}$, with $\gamma^{+i}$ as in Section~\ref{dp:sec:nonambient}, of absolute value at most $\norm z_\tau$, as $\abs{T_\alpha}\le1$ and $\tau\ge1$. So no entry of $h\in\mathcal E_{T,\tau}$ has absolute value above $\norm h_{T,\tau}$. Thus the coordinate forms of $\mathcal E_{T,\tau}$ are bounded and separate points. By Section~\ref{dp:sec:conv}, $\pi$ is a norm on $\mathcal E_{T,\tau}^{\otimes p}$, $\norm\cdot_\Gamma$ is a norm on $\Gamma^p(\mathcal E_{T,\tau})$, and $\pi\le\norm\cdot_\Gamma$.

\emph{Step 2 (the projective norm).} In $\mathcal E_{T,\tau}$ we have $\norm{e_i}\le1$ and $\norm{\partial^\gamma T}\le\tau$, by~\eqref{dp:eq:weighted-norm}. For a word $(i_1,\dots,i_{p-1})\in[m]^{p-1}$ let $\gamma(i_1,\dots,i_{p-1})$ be its content. The value of $T$ at $(i_1,\dots,i_{p-1},i)$ is the $i$-th entry of $\partial^{\gamma(i_1,\dots,i_{p-1})}T$, so
\[
 T=\sum_{(i_1,\dots,i_{p-1})\in[m]^{p-1}}e_{i_1}\otimes\dots\otimes e_{i_{p-1}}\otimes\partial^{\gamma(i_1,\dots,i_{p-1})}T.
\]
This is a sum of $m^{p-1}$ elementary tensors whose factors have norms with product at most $\tau$. Hence $\pi(T)\le\tau m^{p-1}$ in $\mathcal E_{T,\tau}^{\otimes p}$.

\emph{Step 3 (a representation gives a symmetric tensor over $\F_p$).} The elements of $K$ of absolute value below $1$ form an $\F_p$\nobreakdash-linear subspace that does not contain $1$. Extend a basis of this subspace, together with $1$, to a basis of $K$ over $\F_p$. So there is an $\F_p$\nobreakdash-linear map $\rho\colon K\to\F_p$ with $\rho(1)=1$ that vanishes on it. Then $\rho$ fixes $\F_p$, and $\abs x\ge1$ for every $x\in K$ with $\rho(x)\ne0$.

Let $\delta>0$ and $T=\sum_lc_lh_l^{\otimes p}$ with $c_l\in K$, $h_l\in\mathcal E_{T,\tau}$ and $\sum_l\abs{c_l}\norm{h_l}_{T,\tau}^p\le\norm T_\Gamma+\delta$. Choose $\hat h_l\in K^{\mathcal J_m}$ with $q_T(\hat h_l)=h_l$ and $\norm{\hat h_l}_\tau\le(1+\delta)\norm{h_l}_{T,\tau}$, with $\hat h_l=0$ if $h_l=0$. The tensor $\widetilde S=\sum_lc_l\hat h_l^{\otimes p}$ is symmetric, and $q_T^{\otimes p}(\widetilde S)=T$. For $z\in K^{\mathcal J_m}$ the sum of the numbers $\tau^{\abs w_\partial}\abs{z^{\otimes p}(w)}$, $w\in\mathcal J_m^p$, is $\norm z_\tau^p$. So, by the triangle inequality,
\begin{equation}\label{dp:eq:weighted-lift}
 \sum_{w\in\mathcal J_m^p}\tau^{\abs w_\partial}\abs{\widetilde S(w)}\le\sum_l\abs{c_l}\norm{\hat h_l}_\tau^p\le(1+\delta)^p(\norm T_\Gamma+\delta).
\end{equation}
Let $S$ be the tensor over $\F_p$ with the values $S(w)=\rho(\widetilde S(w))$. It is symmetric, as $\widetilde S$ is, so $S\in\Gamma^p(\F_p^{\mathcal J_m})$. As the matrix of $q_T$ has entries in $\F_p$ and $\rho$ is $\F_p$\nobreakdash-linear, each value of $q_T^{\otimes p}(S)$ is the image under $\rho$ of the corresponding value of $q_T^{\otimes p}(\widetilde S)=T$. These values lie in $\F_p$, so $q_T^{\otimes p}(S)=T$. For $\mathbf a\in\operatorname{supp}S$ and a word $w$ of content $\mathbf a$ we have $\rho(\widetilde S(w))=S_{\mathbf a}\ne0$, so $\abs{\widetilde S(w)}\ge1$; and $\abs w_\partial=\abs{\mathbf a}_\partial$. Taking one such word for each $\mathbf a\in\operatorname{supp}S$, \eqref{dp:eq:weighted-lift} gives $\sum_{\mathbf a\in\operatorname{supp}S}\tau^{\abs{\mathbf a}_\partial}\le(1+\delta)^p(\norm T_\Gamma+\delta)$. As $\Gamma^p(\F_p^{\mathcal J_m})$ is finite, one $S$ occurs for arbitrarily small $\delta$. So
\begin{equation}\label{dp:eq:support-bound}
 \norm T_\Gamma\ge\sum_{\mathbf a\in\operatorname{supp}S}\tau^{\abs{\mathbf a}_\partial}\quad\text{for some }S\in\Gamma^p(\F_p^{\mathcal J_m})\text{ with }q_T^{\otimes p}(S)=T.
\end{equation}

\emph{Step 4 (supports with few nonlinear positions).} Let $S\in\Gamma^p(\F_p^{\mathcal J_m})$. Let $\mathcal K_S\subseteq\N^m_p$ be the set of the $\gamma^{+i}$, $i\in[m]$, for the $\gamma\in\N^m_{p-1}$ at which some nonlinear $\mathbf a\in\operatorname{supp}S$ has a nonzero entry. A position has nonzero entries at no more than $p$ elements of $\N^m_{p-1}$, so $\abs{\mathcal K_S}$ is at most $pm$ times the number of nonlinear positions in $\operatorname{supp}S$. We show that at most $p^{\abs{\mathcal K_S}+\binom{m+p-1}p-\mathsf r_m}$ tensors $T'\in\Gamma^p(\F_p^m)$ satisfy $q_{T'}^{\otimes p}(S)=T'$.

Write $S=S^{(0)}+S^{(1)}+S^{(2)}$, where $S^{(0)}$, $S^{(1)}$ and $S^{(2)}$ are supported in the positions $\mathbf a$ with $\abs{\mathbf a}_\partial=0$, $\abs{\mathbf a}_\partial=1$ and $\abs{\mathbf a}_\partial\ge2$. The positions with $\abs{\mathbf a}_\partial=0$ are the $\alpha\in\N^m_p$, extended by $0$, and $q_{T'}$ fixes the $e_i$. So $q_{T'}^{\otimes p}(S^{(0)})=\sum_{\alpha\in\N^m_p}S_\alpha O_\alpha$ does not depend on $T'$.

For $\gamma,\gamma'\in\N^m_{p-1}$ let $\mathbf a(\gamma,\gamma')$ be the position with the entry $1$ at $\gamma$, the entries of $\gamma'$ on $[m]$ and $0$ elsewhere. These are the positions with $\abs{\mathbf a}_\partial=1$. Also $O_{\mathbf a(\gamma,\gamma')}=e_\gamma\cdot O_{\gamma'}$, with $\gamma'$ extended by $0$, since each word of content $\mathbf a(\gamma,\gamma')$ arises from exactly one of the $p$ permutations that define this product, applied to exactly one tensor $e_\gamma\otimes e_{i_1}\otimes\dots\otimes e_{i_{p-1}}$ with $(i_1,\dots,i_{p-1})$ of content $\gamma'$. As the product commutes with $q_{T'}$, and $q_{T'}$ fixes the $e_i$,
\[
 q_{T'}^{\otimes p}(S^{(1)})=\sum_{\gamma,\gamma'\in\N^m_{p-1}}S_{\mathbf a(\gamma,\gamma')}\,(\partial^\gamma T')\cdot O_{\gamma'}=\mathsf L^\sharp(T'),
\]
with $\mathsf L_{\gamma\gamma'}=S_{\mathbf a(\gamma,\gamma')}$ in Lemma~\ref{dp:lem:linear}, for $\Bbbk=\F_p$.

Finally, $q_{T'}^{\otimes p}(S^{(2)})$ depends only on the vectors $q_{T'}(e_\gamma)=\partial^\gamma T'$ for the $\gamma$ at which some $\mathbf a\in\operatorname{supp}S^{(2)}$ has a nonzero entry. That is, it depends only on the coefficients $T'_\alpha$, $\alpha\in\mathcal K_S$. So if $q_{T'}^{\otimes p}(S)=T'$, then $(\id-\mathsf L^\sharp)(T')$ equals $q_{T'}^{\otimes p}(S^{(0)})+q_{T'}^{\otimes p}(S^{(2)})$, which depends only on $(T'_\alpha)_{\alpha\in\mathcal K_S}$. For each of the $p^{\abs{\mathcal K_S}}$ choices of these coefficients the solutions $T'$ lie in one coset of the kernel of $\id-\mathsf L^\sharp$. This kernel has at most $p^{\binom{m+p-1}p-\mathsf r_m}$ elements, by Lemma~\ref{dp:lem:linear}(2), as $\dim\Gamma^p(\F_p^m)=\binom{m+p-1}p$.

\emph{Step 5 (item (1)).} Put $\mathsf n_1=\mathsf r_m/(3p^2(1+\log_pm))$, $\mathsf n_2=\mathsf r_m/(3pm)$ and $\tau=\max\{1,(m/(p(1+\log_pm)))^{1/2}\}$. There are $\binom{\abs{\mathcal J_m}+p-1}p\le\abs{\mathcal J_m}^p$ positions, and $\abs{\mathcal J_m}=m+\binom{m+p-2}{p-1}\le m+m^{p-1}\le2m^{p-1}$. So $p$ times the number of positions is at most $p\,2^pm^{p(p-1)}\le p^{p^2}m^{p^2}=p^{p^2(1+\log_pm)}$, as $p\,2^p\le p^{p^2}$.

Every $S\in\Gamma^p(\F_p^{\mathcal J_m})$ that is nonzero at fewer than $\mathsf n_1$ positions is the sum of $\lceil\mathsf n_1\rceil-1$ tensors $O_{\mathbf a}$, $\mathbf a$ a position, each multiplied by an element of $\F_p$. By the bound on the number of positions above, and as $\lceil\mathsf n_1\rceil-1<\mathsf n_1$, there are at most $p^{\mathsf n_1p^2(1+\log_pm)}=p^{\mathsf r_m/3}$ of them. If moreover fewer than $\mathsf n_2$ of its positions are nonlinear, then $\abs{\mathcal K_S}\le pm\mathsf n_2=\mathsf r_m/3$, and by Step~4 at most $p^{\binom{m+p-1}p-2\mathsf r_m/3}$ tensors $T'$ satisfy $q_{T'}^{\otimes p}(S)=T'$. So at most $p^{\binom{m+p-1}p-\mathsf r_m/3}$ of the $p^{\binom{m+p-1}p}$ elements $T'$ of $\Gamma^p(\F_p^m)$ satisfy $q_{T'}^{\otimes p}(S)=T'$ for such an $S$.

As $\mathsf r_m>0$, these are not all of $\Gamma^p(\F_p^m)$. Choose $T\in\Gamma^p(\F_p^m)$ outside them. Then every $S$ with $q_T^{\otimes p}(S)=T$ is nonzero at no fewer than $\mathsf n_1$ positions, or at no fewer than $\mathsf n_2$ nonlinear positions. The weights $\tau^{\abs{\mathbf a}_\partial}$ in~\eqref{dp:eq:support-bound} are at least $1$, and at least $\tau^2$ for nonlinear $\mathbf a$. Also $\tau^2\mathsf n_2\ge\mathsf n_1$, as $\tau^2\ge m/(p(1+\log_pm))$. So $\norm T_\Gamma\ge\mathsf n_1>0$ in $\Gamma^p(\mathcal E_{T,\tau})$, and $T\ne0$, so $\pi(T)>0$ by Step~1. If $\tau>1$, Step~2 gives
\[
 \frac{\norm T_\Gamma}{\pi(T)}\ge\frac{\mathsf n_1}{\tau m^{p-1}}=\frac{(p-1)!}{3p^{3/2}(2p-1)!}\Bigl(\frac m{1+\log_pm}\Bigr)^{1/2}.
\]
If $\tau=1$, then $m\le p(1+\log_pm)$, so the right side of~\eqref{dp:eq:degree-p-ratio} is at most $(p-1)!\,\pi(T)/(3p(2p-1)!)<\pi(T)\le\norm T_\Gamma$, by Step~1.

\emph{Step 6 (item (2): the law).} Let $T^{(m)},\tau^{(m)},\mathcal E_m$, $m\ge1$, and $\mathcal E$ be as in item~(2), let $n\ge p$, and let $\mathsf g_n\colon K\oplus\mathcal E\to\Gamma^p(\mathcal E)$ be the map $\mathsf g_n(\lambda,h)=\lambda^{n-p}h^{\otimes p}$, with $\lambda^0=1$. Each norm of $\mathcal E_m$ lies between the maximum of the absolute values of the entries, by Step~1, and their sum. So if $K$ is complete, the spaces $\mathcal E_m$ are complete, and so is their $\ell^1$-sum $\mathcal E$. By Step~1 the coordinate forms of the spaces $\mathcal E_m$, composed with the projections of $\mathcal E$ onto them, are bounded, and they separate points. So $\pi$ is a norm on $\mathcal E^{\otimes p}$, and $\norm\cdot_\Gamma\ge\pi$ is a norm on $\Gamma^p(\mathcal E)$.

For finitely many vectors $(\lambda_l,h_l)$ of $K\oplus\mathcal E$ and scalars $s_l$, the value of $\mathsf g_n$ at $\sum_ls_l(\lambda_l,h_l)$ is $(\sum_ls_l\lambda_l)^{n-p}(\sum_ls_lh_l)^{\otimes p}$. By~\eqref{dp:eq:orbit} this is a homogeneous polynomial of degree $n$ in the $s_l$ with coefficients in $\Gamma^p(\mathcal E)$. Moreover $\norm{\lambda^{n-p}h^{\otimes p}}_\Gamma\le\abs\lambda^{n-p}\norm h^p\le\norm{(\lambda,h)}^n$. So $\mathsf g_n$ is a bounded polynomial law of degree $n$ and norm at most $1$.

\emph{Step 7 (item (2): the ratio on $\Gamma^p(\mathcal E)$).} The inclusion of $\mathcal E_m$ in $\mathcal E$ is isometric, and the projection of $\mathcal E$ onto $\mathcal E_m$ is a contractive left inverse. Let $t\in\Gamma^p(\mathcal E_m)$. Then $\pi(t)$ computed in $\mathcal E^{\otimes p}$ is at most $\pi(t)$ computed in $\mathcal E_m^{\otimes p}$. Also $\norm t_\Gamma$ computed in $\Gamma^p(\mathcal E)$ is at least $\norm t_\Gamma$ computed in $\Gamma^p(\mathcal E_m)$, since the $p$-th tensor power of the projection fixes $t$ and does not increase the cost $\sum_l\abs{c_l}\norm{h_l}^p$ of a representation. By item~(1), applied to $t=T^{(m)}\ne0$, and as $m/(1+\log_pm)\to\infty$ with $m$, the quotient $\norm\cdot_\Gamma/\pi$ is unbounded on the nonzero elements of $\Gamma^p(\mathcal E)$.

\emph{Step 8 (item (2): comparing coefficients).} Suppose that $j\mathsf g_n(\lambda,h)=\beta((\lambda,h),\dots,(\lambda,h))$ for all $(\lambda,h)\in K\oplus\mathcal E$, for a linear isometry $j$ of $\Gamma^p(\mathcal E)$, or of its completion, into a normed space $Z$ and a bounded $n$\nobreakdash-linear map $\beta\colon(K\oplus\mathcal E)^n\to Z$. Put $\mathbf 1=(1,0)$ and $h^+=(0,h)$ for $h\in\mathcal E$. For $h_1,\dots,h_p\in\mathcal E$ let $\beta'(h_1,\dots,h_p)$ be the sum, over the $\binom np$ sets of $p$ of the $n$ slots, of the value of $\beta$ at the $n$-tuple with $h_1^+,\dots,h_p^+$ in these slots, in this order, and $\mathbf 1$ in the others. Then $\beta'$ is $p$\nobreakdash-linear, and $\norm{\beta'}\le\binom np\norm\beta$, as $\norm{\mathbf 1}=1$ and $\norm{h^+}=\norm h$.

For $h\in\mathcal E$ and $s=(s_1,s_2)\in K^2$ we have $(s_2,s_1h)=s_2\mathbf 1+s_1h^+$. So $\beta((s_2,s_1h),\dots,(s_2,s_1h))$ is a homogeneous polynomial of degree $n$ in $s$ whose coefficient of $s_1^ps_2^{n-p}$ is $\beta'(h,\dots,h)$. On the other side, $j\mathsf g_n(s_2,s_1h)=s_1^ps_2^{n-p}j(h^{\otimes p})$. As $K$ is infinite, comparing coefficients gives $j(h^{\otimes p})=\beta'(h,\dots,h)$. Let $\beta'_\otimes\colon(\mathcal E^{\otimes p},\pi)\to Z$ be the linear map with $\beta'_\otimes(h_1\otimes\dots\otimes h_p)=\beta'(h_1,\dots,h_p)$, of norm $\norm{\beta'}$. As $\Gamma^p(\mathcal E)$ is spanned by powers, $j(t)=\beta'_\otimes(t)$ for $t\in\Gamma^p(\mathcal E)$. So $\norm t_\Gamma=\norm{j(t)}\le\binom np\norm\beta\,\pi(t)$ for $t\in\Gamma^p(\mathcal E)$, against Step~7. So $\mathsf g_n$ is not ambient of degree $n$, also as a map into the completion of $\Gamma^p(\mathcal E)$.
\end{proof}
